\chapter{Sewing Maps and Periodic Fibers}
\index{sewing map|textbf}
\index{periodic fiber|textbf}
\index{reciprocal sewing|textbf}
\index{Wilson loop!sewing map}

\label{ch:training-sewing-periodic-fibers}

\chapterstatus{foundational mathematical account of periodic-fiber
identification and its representation dependence. The chapter distinguishes
exact maps from finite represented maps and separates reciprocal closure from
stronger physical interpretations such as polarization or topology.}

\section{Learning goals}

After working through this chapter, a reader should be able to

\begin{enumerate}
  \item explain why reciprocal periodicity does not imply literal equality of
        endpoint coordinate arrays;
  \item distinguish a basis inventory, a basis embedding, a Bloch phase
        convention, and a sewing map;
  \item state the domain, codomain, direction, reciprocal shift, and matrix
        convention of a reciprocal sewing map;
  \item derive identity and embedding-phase sewing from declared Bloch sums;
  \item distinguish reciprocal, permutation--phase, plane-wave-index,
        projected-subspace, symmetry, and finite-boundary sewing;
  \item determine from the represented state spaces and conventions which
        sewing construction is applicable;
  \item determine when a projected sewing matrix is unitary;
  \item construct a sewn discrete loop product and prove its gauge covariance;
  \item explain why a finite-cutoff index shift need not be unitary; and
  \item state why full-space holonomy is not automatically evidence of
        topology, polarization, localization, or reduced-model fidelity.
\end{enumerate}

Chapter~\ref{ch:training-periodic-geometry} defined reciprocal-lattice classes
and representatives. Chapter~\ref{ch:training-projection-retained-spaces}
distinguished a retained subspace from a frame used to represent it. Those two
ideas meet here. A reciprocal endpoint can represent the same point of the
Brillouin-zone quotient while belonging to a differently represented fiber.
The operation that identifies those fibers must therefore be explicit.

The subject is narrower than arbitrary basis transformation. Ordinary
reciprocal sewing uses one admitted basis identity at two reciprocal
representatives separated by a reciprocal-lattice vector. If source and target
use different scientific basis identities, the required operation is a basis
transformation, gauge bridge, projection, or other identified-space map. Calling
all such operations ``sewing'' would erase the state spaces being related.

\section{Periodic fibers and correlated basis data}
\index{basis identity!sewing prerequisite}
\index{basis embedding!matrix ordering}

Let reduced reciprocal coordinates be denoted by
\(\mathbf q\in\mathbb R^d\). Integer reciprocal shifts are written
\(\mathbf m\in\mathbb Z^d\), so that \(\mathbf q\) and
\(\mathbf q+\mathbf m\) belong to the same reciprocal-lattice class.

\begin{definition}[Represented periodic-fiber family]
\label{def:training-periodic-fiber-family}
A represented periodic-fiber family for one admitted basis inventory
\(\mathcal B\) is a collection
\begin{equation}
  \left\{\mathcal H_{\mathbf q}^{(\mathcal B)}
  \;\middle|\;\mathbf q\in\mathbb R^d\right\}
\end{equation}
of finite-dimensional complex inner-product spaces, together with an ordered
basis or frame convention at every chosen reciprocal representative. The label
\(\mathcal B\) identifies the correlated scientific basis inventory; it does
not by itself specify geometry, embedding, gauge, energy reference, or a sewing
map.
\end{definition}

Here \(\mathcal B\) denotes a fixed ordered basis inventory rather than a
particular naming scheme. Two labels that happen to look alike do not establish
that the corresponding ordered states, orbital content, spin content, or matrix
indices coincide.

\begin{definition}[Periodic basis embedding]
\label{def:training-periodic-basis-embedding}
For an ordered rank-\(r\) basis inventory \(\mathcal B\), a periodic basis
embedding is the ordered tuple
\begin{equation}
  \mathsf S_{\mathcal B}
  =\left(\mathbf s_1,\ldots,\mathbf s_r\right),
  \qquad
  \mathbf s_a\in\mathbb R^d,
\end{equation}
where \(\mathbf s_a\) is the dimensionless fractional position of matrix index
\(a\) in the ordered primitive direct-lattice basis. Tuple position is the
matrix-index correlation.
\end{definition}

Fractional positions remain as supplied. Reducing every position into a chosen
unit cell can erase an integer lattice lift and thereby change a gauge or
provenance convention. Duplicate positions are valid: several orbitals, spin
states, or other basis states may occupy the same embedded point. The embedding
is not a substitute for the unit-bearing direct-lattice geometry.

\begin{definition}[Positive Bloch basis conventions]
\label{def:training-positive-bloch-conventions}
For localized basis states \(\lvert\mathbf n,a\rangle\), with
\(\mathbf n\in\mathbb Z^d\), the positive cell-phase convention is
\begin{equation}
  \lvert a,\mathbf q\rangle_{\mathrm{cell}}
  =\sum_{\mathbf n}
   e^{2\pi i\mathbf q\cdot\mathbf n}
   \lvert\mathbf n,a\rangle.
  \label{eq:training-cell-phase-bloch-sum}
\end{equation}
The positive position-phase convention is
\begin{equation}
  \lvert a,\mathbf q\rangle_{\mathrm{pos}}
  =\sum_{\mathbf n}
   e^{2\pi i\mathbf q\cdot(\mathbf n+\mathbf s_a)}
   \lvert\mathbf n,a\rangle.
  \label{eq:training-position-phase-bloch-sum}
\end{equation}
The sign, the factor \(2\pi\), and the placement of \(\mathbf s_a\) are parts
of the convention.
\end{definition}

A convention cannot be recovered reliably from a Hamiltonian spectrum. Two
conventions can produce identical eigenvalues while assigning different
endpoint phases to their coordinate bases.

\section{Reciprocal-shift sewing}
\index{reciprocal shift!sewing map}

\begin{definition}[Reciprocal-shift sewing map]
\label{def:training-reciprocal-shift-sewing-map}
For one admitted basis identity \(\mathcal B\), a reciprocal-shift sewing map is
a declared linear map
\begin{equation}
  S_{\mathbf m}(\mathbf q):
  \mathcal H_{\mathbf q}^{(\mathcal B)}
  \longrightarrow
  \mathcal H_{\mathbf q+\mathbf m}^{(\mathcal B)},
  \qquad \mathbf m\in\mathbb Z^d.
  \label{eq:training-reciprocal-sewing-map}
\end{equation}
Its specification includes the source representative, target representative,
reciprocal shift, direction, ordered basis identity, Bloch convention, and
finite matrix representation. It is an \emph{exact unitary reciprocal sewing}
when \(S_{\mathbf m}^{\dagger}S_{\mathbf m}=I\) and the map identifies the
complete represented fibers.
\end{definition}

A coherent exact family additionally satisfies
\begin{equation}
  S_{\mathbf 0}(\mathbf q)=I,
  \qquad
  S_{\mathbf n}(\mathbf q+\mathbf m)
  S_{\mathbf m}(\mathbf q)
  =S_{\mathbf m+\mathbf n}(\mathbf q).
  \label{eq:training-sewing-composition-law}
\end{equation}
A finite represented map need not satisfy exact unitarity or the composition
law after truncation. Such defects characterize the finite representation. A
nearest-unitary replacement is a different mathematical operation and must be
distinguished from the original represented map.

\subsection{Identity reciprocal sewing}

\begin{definition}[Identity reciprocal sewing]
\label{def:training-identity-reciprocal-sewing}
Identity reciprocal sewing is the explicit reciprocal-shift map
\begin{equation}
  S_{\mathbf m}(\mathbf q)=I
\end{equation}
for a basis convention in which the ordered coordinate basis is exactly
periodic under \(\mathbf q\mapsto\mathbf q+\mathbf m\).
\end{definition}

\begin{proposition}[Cell-phase reciprocal periodicity]
\label{prop:training-cell-phase-identity-sewing}
For the Bloch sums in
Equation~\eqref{eq:training-cell-phase-bloch-sum}, every integer reciprocal
shift has identity sewing.
\end{proposition}

\begin{proof}
For \(\mathbf m,\mathbf n\in\mathbb Z^d\),
\begin{equation}
  e^{2\pi i(\mathbf q+\mathbf m)\cdot\mathbf n}
  =e^{2\pi i\mathbf q\cdot\mathbf n}
   e^{2\pi i\mathbf m\cdot\mathbf n}
  =e^{2\pi i\mathbf q\cdot\mathbf n}.
\end{equation}
Thus every ordered basis vector at \(\mathbf q+\mathbf m\) equals its correlated
basis vector at \(\mathbf q\) in this convention.
\end{proof}

Identity is therefore a result of declared basis periodicity, not a default
inference from periodic eigenvalues.

\subsection{Embedding-position reciprocal sewing}

\begin{definition}[Embedding-position reciprocal sewing]
\label{def:training-embedding-position-sewing}
For the positive position-phase convention and embedding
\(\mathsf S_{\mathcal B}\), embedding-position reciprocal sewing is the
diagonal map
\begin{equation}
  S_{\mathbf m}^{\mathrm{pos}}
  =\operatorname{diag}\left(
    e^{2\pi i\mathbf m\cdot\mathbf s_1},\ldots,
    e^{2\pi i\mathbf m\cdot\mathbf s_r}
  \right),
  \label{eq:training-embedding-position-sewing}
\end{equation}
with matrix order inherited from the basis inventory.
\end{definition}

\begin{proposition}[Embedding sewing is a unitary representation]
\label{prop:training-embedding-sewing-unitary}
The maps in Equation~\eqref{eq:training-embedding-position-sewing} are unitary
and satisfy
\begin{equation}
  S_{\mathbf n}^{\mathrm{pos}}
  S_{\mathbf m}^{\mathrm{pos}}
  =S_{\mathbf m+\mathbf n}^{\mathrm{pos}}.
\end{equation}
\end{proposition}

\begin{proof}
Every diagonal entry has modulus one, so the matrix is unitary. Multiplying two
entries gives
\begin{equation}
  e^{2\pi i\mathbf n\cdot\mathbf s_a}
  e^{2\pi i\mathbf m\cdot\mathbf s_a}
  =e^{2\pi i(\mathbf m+\mathbf n)\cdot\mathbf s_a},
\end{equation}
which proves the composition law index by index.
\end{proof}

The sign in Equation~\eqref{eq:training-embedding-position-sewing} follows from
the positive phase in
Equation~\eqref{eq:training-position-phase-bloch-sum}. A negative Bloch phase
would require the conjugate convention. Copying the matrix while changing the
Bloch-sum sign would be a convention error.

\subsection{Permutation--phase reciprocal sewing}

\begin{definition}[Permutation--phase reciprocal sewing]
\label{def:training-permutation-phase-sewing}
Suppose a reviewed endpoint convention relates the same ordered scientific
basis inventory by a permutation \(\pi\) and phases \(\phi_a\). The
permutation--phase sewing map is
\begin{equation}
  S_{\pi,\boldsymbol\phi}\mathbf e_a
  =e^{i\phi_a}\mathbf e_{\pi(a)},
  \qquad a=1,\ldots,r.
  \label{eq:training-permutation-phase-sewing}
\end{equation}
The permutation and phases are declared correlation data, not inferred repairs
for an accidental ordering mismatch.
\end{definition}

This map is unitary because it is a permutation matrix multiplied by unit-modulus
phases. It does not authorize matching states by nearest eigenvalue, largest
overlap, orbital-name similarity, or post-hoc sorting. Those are separate
correspondence procedures with their own ambiguity and evidence boundaries.

\section{Plane-wave index sewing}
\index{plane-wave basis!reciprocal sewing}
\index{finite cutoff!nonunitary sewing}

A plane-wave fiber at \(\mathbf q\) commonly stores coefficients indexed by
reciprocal vectors \(\mathbf G\). The physical wavevector represented by an
index is \(\mathbf q+\mathbf G\). Under
\(\mathbf q\mapsto\mathbf q+\mathbf m\), preservation of the same physical
wavevector requires
\begin{equation}
  \mathbf G' = \mathbf G-\mathbf m.
  \label{eq:training-plane-wave-index-shift}
\end{equation}

\begin{definition}[Plane-wave index sewing]
\label{def:training-plane-wave-index-sewing}
Let \(\mathcal I_{\mathbf q}\) and
\(\mathcal I_{\mathbf q+\mathbf m}\) be the finite reciprocal-index inventories
at two endpoints. The finite plane-wave index-sewing map is
\begin{equation}
  P_{\mathbf m}\mathbf e_{\mathbf G}
  =
  \begin{cases}
    \mathbf e_{\mathbf G-\mathbf m},
      &\mathbf G-\mathbf m\in\mathcal I_{\mathbf q+\mathbf m},\\
    0,
      &\mathbf G-\mathbf m\notin\mathcal I_{\mathbf q+\mathbf m}.
  \end{cases}
  \label{eq:training-finite-plane-wave-sewing}
\end{equation}
The source and target inventories, index ordering, reciprocal shift, and omitted
indices are parts of the represented map.
\end{definition}

On a complete reciprocal-index lattice this is a permutation and hence unitary.
At finite cutoff, an index can leave the target inventory or a target index can
lack a source preimage. The represented map can then be rectangular, rank
deficient, or merely a partial isometry.

\begin{proposition}[Finite-cutoff unitarity criterion]
\label{prop:training-plane-wave-sewing-unitarity}
The finite plane-wave map in
Equation~\eqref{eq:training-finite-plane-wave-sewing} is unitary exactly when
\(\mathbf G\mapsto\mathbf G-\mathbf m\) is a bijection from
\(\mathcal I_{\mathbf q}\) onto
\(\mathcal I_{\mathbf q+\mathbf m}\).
\end{proposition}

\begin{proof}
If the index shift is a bijection, the matrix has exactly one unit entry in every
row and column and is a permutation matrix. Conversely, failure of injectivity
or surjectivity prevents the columns or rows from forming a complete
orthonormal basis, so the matrix cannot be unitary.
\end{proof}

A nearest-unitary projection would answer a different numerical question. It
must not be substituted silently for the finite-cutoff map or presented as
proof that the cutoff representation closes exactly.

\section{Projected-subspace sewing}
\index{projected sewing|textbf}
\index{retained subspace!sewing}

Let
\begin{equation}
  F_{\mathrm s}:\mathbb C^r\longrightarrow
  \mathcal H_{\mathbf q}^{(\mathcal B)},
  \qquad
  F_{\mathrm t}:\mathbb C^r\longrightarrow
  \mathcal H_{\mathbf q+\mathbf m}^{(\mathcal B)}
\end{equation}
be isometric frames, so
\(F_{\mathrm s}^{\dagger}F_{\mathrm s}
=F_{\mathrm t}^{\dagger}F_{\mathrm t}=I_r\).

\begin{definition}[Projected-subspace sewing matrix]
\label{def:training-projected-subspace-sewing}
Given an ambient sewing map \(S_{\mathbf m}\), the projected-subspace sewing
matrix from the source retained coordinates to the target retained coordinates
is
\begin{equation}
  C_{\mathbf m}
  =F_{\mathrm t}^{\dagger}
   S_{\mathbf m}
   F_{\mathrm s}.
  \label{eq:training-projected-sewing}
\end{equation}
The selected source and target subspaces are part of the definition. Projection
does not itself select bands or infer an occupied space.
\end{definition}

\begin{proposition}[Criterion for unitary projected sewing]
\label{prop:training-projected-sewing-unitary}
Assume \(S_{\mathbf m}\) is unitary and the source and target frames have equal
rank. Then \(C_{\mathbf m}\) is unitary if and only if
\begin{equation}
  S_{\mathbf m}\operatorname{Ran}F_{\mathrm s}
  =\operatorname{Ran}F_{\mathrm t}.
  \label{eq:training-sewn-subspace-equality}
\end{equation}
\end{proposition}

\begin{proof}
Let \(P_{\mathrm t}=F_{\mathrm t}F_{\mathrm t}^{\dagger}\). Then
\begin{equation}
  C_{\mathbf m}^{\dagger}C_{\mathbf m}
  =F_{\mathrm s}^{\dagger}S_{\mathbf m}^{\dagger}
   P_{\mathrm t}S_{\mathbf m}F_{\mathrm s}.
\end{equation}
This equals \(I_r\) exactly when every vector in
\(S_{\mathbf m}\operatorname{Ran}F_{\mathrm s}\) lies in
\(\operatorname{Ran}F_{\mathrm t}\). Equal finite ranks then promote the
inclusion to equality. The reverse product gives the same condition from the
target side.
\end{proof}

Singular values below one therefore have scientific meaning: they can indicate
subspace mismatch, an incomplete ambient map, finite precision, or a combination
of these. A polar factor can define a nearest unitary alignment when that
operation is requested, but the polar factor is not the original projected
sewing matrix.

\section{Symmetry sewing}
\index{symmetry sewing}
\index{antiunitary symmetry!sewing}

\begin{definition}[Symmetry sewing transformation]
\label{def:training-symmetry-sewing}
For a declared spatial or internal symmetry \(g\), a symmetry sewing
transformation relates the fibers at symmetry-connected reciprocal points,
\begin{equation}
  U_g(\mathbf q):
  \mathcal H_{\mathbf q}
  \longrightarrow
  \mathcal H_{g\mathbf q}.
\end{equation}
It is linear and unitary for a unitary symmetry, and conjugate-linear and
norm-preserving for an antiunitary symmetry. The symmetry action on reciprocal
coordinates, represented basis states, spin or internal degrees of freedom,
and any reciprocal shift used to return to a chosen representative are all
parts of the transformation.
\end{definition}

Symmetry sewing is not ordinary reciprocal sewing. Even when
\(g\mathbf q=\mathbf q+\mathbf m\), the map can contain a nontrivial symmetry
action in addition to reciprocal closure. Antiunitary sewing cannot be inserted
into a product of complex-linear matrices without explicitly representing the
associated complex conjugation.

\section{Finite-boundary sewing}
\index{finite boundary!sewing}
\index{boundary twist!seam identification}

\begin{definition}[Finite-boundary sewing]
\label{def:training-finite-boundary-sewing}
Finite-boundary sewing is the declared identification of degrees of freedom on
opposite sides of a finite real-space fundamental domain. In one scalar
direction a twist boundary condition has the form
\begin{equation}
  \psi(x+L)=e^{i\vartheta}\psi(x),
  \label{eq:training-finite-boundary-twist}
\end{equation}
with analogous matrix-valued identifications for multicomponent states. The
boundary orientation, paired boundary degrees of freedom, twist, and represented
state space are parts of the sewing data.
\end{definition}

Finite-boundary sewing acts at a real-space quotient seam. Reciprocal sewing
identifies fibers at reciprocal representatives. Fourier duality can relate the
two constructions, but they are not the same map and should not be conflated
under one undifferentiated notion of sewing.

\section{Applicability conditions for the sewing constructions}
\index{sewing map!applicability conditions}

\begin{definition}[Applicability of a sewing construction]
\label{def:training-sewing-applicability}
A sewing construction is applicable when its source and target state spaces,
the relation between their base points, the ordered basis correlation, and the
phase or symmetry convention agree with the hypotheses of that construction.
The construction is selected from these mathematical data, not from the desired
appearance of the resulting matrix.
\end{definition}

The seven constructions are not seven interchangeable algorithms for the same
problem. They answer different identification questions. Their applicability
conditions are as follows.

\begin{description}
  \item[Identity reciprocal sewing.]
  Identity sewing applies when the source and target are reciprocal representatives
  of the same fiber, the same ordered basis inventory is used at both points,
  and each target coordinate basis vector is literally the correlated source
  basis vector. The cell-phase Bloch convention satisfies this condition for
  every integer reciprocal shift. Other conventions may satisfy it for a
  particular shift when all associated phases are one. Spectral periodicity
  alone is insufficient.

  \item[Embedding-position reciprocal sewing.]
  Embedding-position sewing applies when the represented states are localized at
  known fractional positions, the basis order is unchanged, and the Bloch sum
  includes those positions in its phase. The reciprocal shift then produces
  the diagonal phases in
  Equation~\eqref{eq:training-embedding-position-sewing}. This construction is
  inappropriate when the positions are unknown, when the basis is not
  correlated by matrix index, or when the phase convention omits the embedded
  positions.

  \item[Permutation--phase reciprocal sewing.]
  Permutation--phase sewing applies when source and target represent the same
  scientific basis states but the endpoint convention deliberately reorders
  them or multiplies them by known phases. The correspondence must be known
  independently of the matrix one wishes to repair. It is not justified by
  nearest-energy sorting, largest-overlap matching, or a visually convenient
  relabeling.

  \item[Plane-wave index sewing.]
  Plane-wave index sewing applies when the fiber coordinates are reciprocal-index
  coefficients and equality of physical wavevectors requires the relabeling
  \(\mathbf G\mapsto\mathbf G-\mathbf m\). For a complete reciprocal-index
  lattice this is an exact permutation. For a finite cutoff it is the
  appropriate represented shift only after the source and target inventories
  are stated; it is an exact unitary sewing only when the index shift is a
  bijection.

  \item[Projected-subspace sewing.]
  Projected-subspace sewing applies only after an ambient sewing map and explicit
  source and target retained subspaces have been specified. It is the induced
  map between retained coordinates, not an alternative way to construct the
  ambient identification. It is unitary precisely under the condition in
  Proposition~\ref{prop:training-projected-sewing-unitary}.

  \item[Symmetry sewing.]
  Symmetry sewing applies when the base points are related by a declared spatial,
  internal, or antiunitary symmetry and the identification must include that
  symmetry action. If the symmetry-related point lies outside the chosen
  reciprocal cell, a reciprocal-shift sewing may also be required to return to
  the chosen representative. A unitary spatial symmetry and an antiunitary
  symmetry require different algebraic treatments.

  \item[Finite-boundary sewing.]
  Finite-boundary sewing applies when opposite boundaries of a finite real-space
  domain are identified, possibly with a scalar or matrix-valued twist. It is
  the relevant construction for quotient seams and twisted finite-domain
  boundary conditions. Open, Dirichlet, or Neumann boundaries do not identify
  opposite boundary fibers and therefore do not call for this sewing map. It is
  not the endpoint map for a reciprocal-space loop merely because both problems
  use a phase.
\end{description}

The constructions can be composed when more than one condition is present. A
position-phase basis may also undergo a known endpoint permutation, in which
case the diagonal embedding phase and the permutation--phase map both appear in
a stated order. A symmetry-related point may require a reciprocal shift before
comparison with the chosen representative. Projected-subspace sewing is then
formed from the resulting ambient composition. Because matrix products need not
commute, the direction and order of these maps are part of the mathematics.

\section{Reciprocal loops and sewn loop products}
\index{reciprocal loop!orientation}
\index{Wilson loop!discrete product}
\index{holonomy!gauge covariance}

\begin{definition}[Oriented reciprocal loop]
\label{def:training-oriented-reciprocal-loop}
An oriented reciprocal loop is an ordered sequence
\begin{equation}
  \Gamma=(\mathbf q_0,\mathbf q_1,\ldots,\mathbf q_N;\mathbf m)
\end{equation}
with
\begin{equation}
  \mathbf q_N=\mathbf q_0+\mathbf m,
  \qquad \mathbf m\in\mathbb Z^d.
\end{equation}
Point order fixes orientation, and \(\mathbf m\) is the explicit closing shift.
A half-open storage convention may omit \(\mathbf q_N\), but it must retain
sufficient information to reconstruct the endpoint and its sewing map.
\end{definition}

Let \(F_j:\mathbb C^r\to\mathcal H_{\mathbf q_j}\) be isometric selected
frames. Neighboring fibers require declared comparison maps
\begin{equation}
  J_j:\mathcal H_{\mathbf q_j}
  \longrightarrow\mathcal H_{\mathbf q_{j+1}}.
\end{equation}
A common ambient representation can make \(J_j\) the represented identity, but
that fact should be stated rather than hidden.

\begin{definition}[Sewn discrete loop product]
\label{def:training-sewn-discrete-loop-product}
Define neighbor links and endpoint sewing by
\begin{equation}
  L_j=F_{j+1}^{\dagger}J_jF_j,
  \qquad j=0,\ldots,N-1,
  \label{eq:training-neighbor-link}
\end{equation}
\begin{equation}
  C=F_N^{\dagger}S_{\mathbf m}(\mathbf q_0)F_0.
  \label{eq:training-endpoint-sewing-link}
\end{equation}
The sewn discrete loop product acting on the initial retained coordinates is
\begin{equation}
  W_{\Gamma}
  =C^{\dagger}L_{N-1}\cdots L_1L_0.
  \label{eq:training-sewn-loop-product}
\end{equation}
When the links implement an admitted parallel-transport or Wilson-loop
discretization, \(W_{\Gamma}\) is the corresponding finite Wilson or holonomy
diagnostic. The link construction and any polar unitarization remain part of
the method.
\end{definition}

The factor \(C^{\dagger}\) returns endpoint coordinates to the initial fiber.
Using \(C\) instead reverses this orientation. A phase sign or fractional root
is meaningless unless the map direction is fixed.

\begin{proposition}[Gauge covariance of the sewn loop product]
\label{prop:training-sewn-loop-gauge-covariance}
Replace each selected frame by
\begin{equation}
  F_j'=F_jG_j,
  \qquad G_j\in U(r).
\end{equation}
Then
\begin{equation}
  W_{\Gamma}'=G_0^{\dagger}W_{\Gamma}G_0.
  \label{eq:training-sewn-loop-covariance}
\end{equation}
Consequently, the eigenvalue multiset of \(W_{\Gamma}\) is invariant under
these frame-coordinate changes.
\end{proposition}

\begin{proof}
The links transform as
\begin{equation}
  L_j'=G_{j+1}^{\dagger}L_jG_j,
  \qquad
  C'=G_N^{\dagger}CG_0.
\end{equation}
Substitution into Equation~\eqref{eq:training-sewn-loop-product} cancels every
interior factor \(G_jG_j^{\dagger}\), including the endpoint factor, leaving
\(G_0^{\dagger}W_{\Gamma}G_0\).
\end{proof}

This covariance result does not prove that a particular finite product is
unitary, converged, topological, or physically interpreted correctly. Finite
neighbor overlaps can be nonunitary. Replacing each full-rank link by its polar
factor is a defined numerical construction, not an algebraic identity. Standard
discrete Berry-phase and Wilson-loop methods provide important precedents for
such products~\cite{fukui2005,soluyanov2011}, while Wannier localization supplies
a related gauge-selection setting~\cite{marzari1997,marzari2012}.

\subsection{The full-space limitation}

Suppose a rank-\(r\) frame spans an entire rank-\(r\) ambient model at every
point. The projectors then contain no nontrivial selected-versus-discarded
subspace information. A sewn loop product can still depend on the chosen
embedding, Bloch convention, basis transport, and numerical link prescription.
It must not be interpreted automatically as evidence of nontrivial topology, polarization,
Wannier centers, or fidelity to a larger parent model.

A stronger interpretation requires additional scientific structure, such as an
isolated occupied subspace inside a larger parent space, a retained subspace
selected by an explicit rule, a symmetry-protected invariant with its
hypotheses checked, or a converged family of loops supporting the claimed
quantity.

\section{Numerical diagnostics and admissibility}
\index{sewing defect}
\index{unitarity defect!sewing map}

\begin{definition}[Normalized sewing unitarity defect]
\label{def:training-sewing-unitarity-defect}
For a square rank-\(r\) represented sewing matrix \(S\), the normalized
Frobenius unitarity defect is
\begin{equation}
  d_{\mathrm{unitary}}(S)
  =\frac{\lVert S^{\dagger}S-I_r\rVert_{\mathrm F}}{\sqrt r}.
  \label{eq:training-sewing-unitarity-defect}
\end{equation}
The quantity is meaningful only together with its norm, normalization, matrix
direction, and basis correlation.
\end{definition}

\begin{definition}[Sewn-subspace mismatch]
\label{def:training-sewn-subspace-mismatch}
For equal-rank source and target projectors
\(P_{\mathrm s}=F_{\mathrm s}F_{\mathrm s}^{\dagger}\) and
\(P_{\mathrm t}=F_{\mathrm t}F_{\mathrm t}^{\dagger}\), and a unitary ambient
sewing map \(S\), a normalized projector mismatch is
\begin{equation}
  d_{\mathrm{subspace}}(S;P_{\mathrm s},P_{\mathrm t})
  =\frac{\lVert P_{\mathrm t}-SP_{\mathrm s}S^{\dagger}\rVert_{\mathrm F}}
  {\sqrt r}.
  \label{eq:training-sewn-subspace-mismatch}
\end{equation}
\end{definition}

The two defects answer different questions. The first tests the represented
ambient map. The second tests whether the selected subspaces correspond under
that map. Neither is a band-energy error, hopping-locality measure, or
scientific acceptance threshold.

When a construction requires a unique unitary polar factor, a rank-deficient
projected link is inadmissible for that construction. Singular vectors may still
be computed, but the unitary action on null directions is not uniquely
determined. One must therefore distinguish a full-rank link, a rank-deficient
link under a stated threshold, and a link whose numerical rank has not been
resolved. None of these distinctions justifies replacing the original map by
identity sewing or by an unreported nearest-unitary map.

\section{Relations among the sewing constructions}

The preceding definitions separate several mathematical roles. Periodic
geometry supplies direct and reciprocal bases and reciprocal-lattice classes.
The basis inventory specifies the ordered represented states, while the basis
embedding assigns their fractional positions. A Bloch convention fixes the
phase used to represent those states. Reciprocal sewing identifies two
reciprocal representatives of the same basis; projected sewing restricts that
identification to chosen subspaces. Symmetry sewing adds a spatial or internal
symmetry action, and finite-boundary sewing closes a real-space quotient seam.
A reciprocal loop then combines its ordered links with an explicit endpoint
identification. These constructions can interact, but none is a synonym for
another.

Chapter~\ref{ch:training-isolated-band-fourier-reduction} next studies scalar
Fourier reduction, for which phase conventions and reciprocal closure remain
important even though a scalar energy is gauge invariant. Chapter~\ref{ch:training-composite-band-gauge-alignment}
then applies ambient sewing to frame transport, endpoint closure, distributed
holonomy, and block-hopping locality. Appendices~\ref{app:one-dimensional-reduction}
and~\ref{app:two-dimensional-reduction} provide controlled computational
examples. Those examples inherit the definitions here but retain their own
numerical methods, evidence status, and interpretation limits.
